% Options for packages loaded elsewhere
\PassOptionsToPackage{unicode}{hyperref}
\PassOptionsToPackage{hyphens}{url}
\documentclass[
]{article}
\usepackage{xcolor}
\usepackage[margin=1in]{geometry}
\usepackage{amsmath,amssymb}
\setcounter{secnumdepth}{-\maxdimen} % remove section numbering
\usepackage{iftex}
\ifPDFTeX
  \usepackage[T1]{fontenc}
  \usepackage[utf8]{inputenc}
  \usepackage{textcomp} % provide euro and other symbols
\else % if luatex or xetex
  \usepackage{unicode-math} % this also loads fontspec
  \defaultfontfeatures{Scale=MatchLowercase}
  \defaultfontfeatures[\rmfamily]{Ligatures=TeX,Scale=1}
\fi
\usepackage{lmodern}
\ifPDFTeX\else
  % xetex/luatex font selection
\fi
% Use upquote if available, for straight quotes in verbatim environments
\IfFileExists{upquote.sty}{\usepackage{upquote}}{}
\IfFileExists{microtype.sty}{% use microtype if available
  \usepackage[]{microtype}
  \UseMicrotypeSet[protrusion]{basicmath} % disable protrusion for tt fonts
}{}
\makeatletter
\@ifundefined{KOMAClassName}{% if non-KOMA class
  \IfFileExists{parskip.sty}{%
    \usepackage{parskip}
  }{% else
    \setlength{\parindent}{0pt}
    \setlength{\parskip}{6pt plus 2pt minus 1pt}}
}{% if KOMA class
  \KOMAoptions{parskip=half}}
\makeatother
\usepackage{color}
\usepackage{fancyvrb}
\newcommand{\VerbBar}{|}
\newcommand{\VERB}{\Verb[commandchars=\\\{\}]}
\DefineVerbatimEnvironment{Highlighting}{Verbatim}{commandchars=\\\{\}}
% Add ',fontsize=\small' for more characters per line
\usepackage{framed}
\definecolor{shadecolor}{RGB}{248,248,248}
\newenvironment{Shaded}{\begin{snugshade}}{\end{snugshade}}
\newcommand{\AlertTok}[1]{\textcolor[rgb]{0.94,0.16,0.16}{#1}}
\newcommand{\AnnotationTok}[1]{\textcolor[rgb]{0.56,0.35,0.01}{\textbf{\textit{#1}}}}
\newcommand{\AttributeTok}[1]{\textcolor[rgb]{0.13,0.29,0.53}{#1}}
\newcommand{\BaseNTok}[1]{\textcolor[rgb]{0.00,0.00,0.81}{#1}}
\newcommand{\BuiltInTok}[1]{#1}
\newcommand{\CharTok}[1]{\textcolor[rgb]{0.31,0.60,0.02}{#1}}
\newcommand{\CommentTok}[1]{\textcolor[rgb]{0.56,0.35,0.01}{\textit{#1}}}
\newcommand{\CommentVarTok}[1]{\textcolor[rgb]{0.56,0.35,0.01}{\textbf{\textit{#1}}}}
\newcommand{\ConstantTok}[1]{\textcolor[rgb]{0.56,0.35,0.01}{#1}}
\newcommand{\ControlFlowTok}[1]{\textcolor[rgb]{0.13,0.29,0.53}{\textbf{#1}}}
\newcommand{\DataTypeTok}[1]{\textcolor[rgb]{0.13,0.29,0.53}{#1}}
\newcommand{\DecValTok}[1]{\textcolor[rgb]{0.00,0.00,0.81}{#1}}
\newcommand{\DocumentationTok}[1]{\textcolor[rgb]{0.56,0.35,0.01}{\textbf{\textit{#1}}}}
\newcommand{\ErrorTok}[1]{\textcolor[rgb]{0.64,0.00,0.00}{\textbf{#1}}}
\newcommand{\ExtensionTok}[1]{#1}
\newcommand{\FloatTok}[1]{\textcolor[rgb]{0.00,0.00,0.81}{#1}}
\newcommand{\FunctionTok}[1]{\textcolor[rgb]{0.13,0.29,0.53}{\textbf{#1}}}
\newcommand{\ImportTok}[1]{#1}
\newcommand{\InformationTok}[1]{\textcolor[rgb]{0.56,0.35,0.01}{\textbf{\textit{#1}}}}
\newcommand{\KeywordTok}[1]{\textcolor[rgb]{0.13,0.29,0.53}{\textbf{#1}}}
\newcommand{\NormalTok}[1]{#1}
\newcommand{\OperatorTok}[1]{\textcolor[rgb]{0.81,0.36,0.00}{\textbf{#1}}}
\newcommand{\OtherTok}[1]{\textcolor[rgb]{0.56,0.35,0.01}{#1}}
\newcommand{\PreprocessorTok}[1]{\textcolor[rgb]{0.56,0.35,0.01}{\textit{#1}}}
\newcommand{\RegionMarkerTok}[1]{#1}
\newcommand{\SpecialCharTok}[1]{\textcolor[rgb]{0.81,0.36,0.00}{\textbf{#1}}}
\newcommand{\SpecialStringTok}[1]{\textcolor[rgb]{0.31,0.60,0.02}{#1}}
\newcommand{\StringTok}[1]{\textcolor[rgb]{0.31,0.60,0.02}{#1}}
\newcommand{\VariableTok}[1]{\textcolor[rgb]{0.00,0.00,0.00}{#1}}
\newcommand{\VerbatimStringTok}[1]{\textcolor[rgb]{0.31,0.60,0.02}{#1}}
\newcommand{\WarningTok}[1]{\textcolor[rgb]{0.56,0.35,0.01}{\textbf{\textit{#1}}}}
\usepackage{graphicx}
\makeatletter
\newsavebox\pandoc@box
\newcommand*\pandocbounded[1]{% scales image to fit in text height/width
  \sbox\pandoc@box{#1}%
  \Gscale@div\@tempa{\textheight}{\dimexpr\ht\pandoc@box+\dp\pandoc@box\relax}%
  \Gscale@div\@tempb{\linewidth}{\wd\pandoc@box}%
  \ifdim\@tempb\p@<\@tempa\p@\let\@tempa\@tempb\fi% select the smaller of both
  \ifdim\@tempa\p@<\p@\scalebox{\@tempa}{\usebox\pandoc@box}%
  \else\usebox{\pandoc@box}%
  \fi%
}
% Set default figure placement to htbp
\def\fps@figure{htbp}
\makeatother
\setlength{\emergencystretch}{3em} % prevent overfull lines
\providecommand{\tightlist}{%
  \setlength{\itemsep}{0pt}\setlength{\parskip}{0pt}}
\usepackage{bookmark}
\IfFileExists{xurl.sty}{\usepackage{xurl}}{} % add URL line breaks if available
\urlstyle{same}
\hypersetup{
  pdftitle={Statistical Modeling and Graphing},
  pdfauthor={Ravina Gaikwad},
  hidelinks,
  pdfcreator={LaTeX via pandoc}}

\title{Statistical Modeling and Graphing}
\author{Ravina Gaikwad}
\date{2025-12-20}

\begin{document}
\maketitle

\textbf{Machine:} HQSML-FVFGF5W8Q05N - Darwin 24.6.0

\begin{center}\includegraphics{05-Lab-Solution_files/figure-latex/load-and-explore-data-1} \end{center}

\emph{Citation}: Trends in Internet-based business-to-business marketing

\emph{Abstract}: The Internet is changing the transactional paradigms
under which businesses-to-business marketers operate.
Business-to-business marketers that take advantage of the operational
efficiencies and effectiveness that emerge from utilizing the Internet
in transactions are out performing firms that utilize traditional
transactional processes. As an example, Dell computers, by utilizing
business-to-business processes that take advantage of the Internet, has
gained the largest market share in the PC business when compared to
traditional manufacturers such as Compaq. This paper first examines the
genesis of the Internet movement in business-to-business markets. The
long-term impact of the increase of business-to-business utilization of
the Internet on the marketing theory and marketing process is then
discussed. Finally, managerial implications and directions for future
research are highlighted.

Dataset includes:

\begin{verbatim}
1)  Business marketing focus - traditional or forward thinking.

2)  Internet use - low, medium, or high levels of business marketing use on the internet.

3)  Time _ 1 - sales scores at the first measurement time.

4)  Time _ 2 - sales scores at the second measurement time
\end{verbatim}

On all of these questions, be sure to include a coherent label for the X
and Y axes. You should change them to be ``professional looking''
(i.e.~Proper Case, explain the variable listed, and could be printed in
a journal). The following will be assessed:

\begin{verbatim}
1)  Is it readable?

2)  Is X-axis labeled appropriately?

3)  Is Y-axis labeled appropriately?

4)  Is it the right graph?

5)  Do the labels in the legend look appropriate?

6)  Are there error bars when appropriate?
\end{verbatim}

We won't grade for color of bars or background color, but you should
consider that these things are usually printed in black/white - so be
sure you know how to change those values as well as get rid of that grey
background.

Please note that each subpoint (i.e.~a, b) indicates a different chart.

\begin{enumerate}
\def\labelenumi{\arabic{enumi})}
\item
  Make a simple histogram using ggplot:

  \begin{enumerate}
  \def\labelenumii{\alph{enumii}.}
  \tightlist
  \item
    Sales at time 1
  \end{enumerate}
\end{enumerate}

\begin{Shaded}
\begin{Highlighting}[]
\FunctionTok{ggplot}\NormalTok{(labdata, }\FunctionTok{aes}\NormalTok{(}\AttributeTok{x =}\NormalTok{ time}\FloatTok{.1}\NormalTok{)) }\SpecialCharTok{+}
  \FunctionTok{geom\_histogram}\NormalTok{(}\AttributeTok{binwidth =} \FloatTok{0.5}\NormalTok{, }\AttributeTok{color =} \StringTok{"black"}\NormalTok{, }\AttributeTok{fill =} \StringTok{"gray70"}\NormalTok{) }\SpecialCharTok{+}
  \FunctionTok{labs}\NormalTok{(}\AttributeTok{x =} \StringTok{"Sales Score at Time 1"}\NormalTok{, }
       \AttributeTok{y =} \StringTok{"Frequency"}\NormalTok{,}
       \AttributeTok{title =} \StringTok{"Distribution of Sales Scores at Time 1"}\NormalTok{) }\SpecialCharTok{+}
\NormalTok{  cleanup}
\end{Highlighting}
\end{Shaded}

\pandocbounded{\includegraphics[keepaspectratio]{05-Lab-Solution_files/figure-latex/hist1-1.pdf}}

The above histogram (1a) shows the frequency distribution of sales
scores at Time 1, this shows the overall pattern \& spread of initial
sales performance across all companies.

\begin{verbatim}
b.  Sales at time 2
\end{verbatim}

\begin{Shaded}
\begin{Highlighting}[]
\FunctionTok{ggplot}\NormalTok{(labdata, }\FunctionTok{aes}\NormalTok{(}\AttributeTok{x =}\NormalTok{ time}\FloatTok{.2}\NormalTok{)) }\SpecialCharTok{+}
  \FunctionTok{geom\_histogram}\NormalTok{(}\AttributeTok{binwidth =} \FloatTok{0.5}\NormalTok{, }\AttributeTok{color =} \StringTok{"black"}\NormalTok{, }\AttributeTok{fill =} \StringTok{"gray70"}\NormalTok{) }\SpecialCharTok{+}
  \FunctionTok{labs}\NormalTok{(}\AttributeTok{x =} \StringTok{"Sales Score at Time 2"}\NormalTok{, }
       \AttributeTok{y =} \StringTok{"Frequency"}\NormalTok{,}
       \AttributeTok{title =} \StringTok{"Distribution of Sales Scores at Time 2"}\NormalTok{) }\SpecialCharTok{+}
\NormalTok{  cleanup}
\end{Highlighting}
\end{Shaded}

\pandocbounded{\includegraphics[keepaspectratio]{05-Lab-Solution_files/figure-latex/hist2-1.pdf}}

The above histogram (1b) shows the frequency distribution of sales
scores at Time 2, showing how sales performance changed in the follow-up
measurement period.

\begin{enumerate}
\def\labelenumi{\arabic{enumi})}
\setcounter{enumi}{1}
\item
  Make a bar chart with two independent variables:

  \begin{enumerate}
  \def\labelenumii{\alph{enumii}.}
  \tightlist
  \item
    Business focus, internet, DV: sales at time 2
  \end{enumerate}
\end{enumerate}

\begin{Shaded}
\begin{Highlighting}[]
\CommentTok{\# Calculate means and standard errors}
\NormalTok{summary\_data }\OtherTok{\textless{}{-}}\NormalTok{ labdata }\SpecialCharTok{\%\textgreater{}\%}
  \FunctionTok{group\_by}\NormalTok{(biz\_focus, internet) }\SpecialCharTok{\%\textgreater{}\%}
  \FunctionTok{summarise}\NormalTok{(}
    \AttributeTok{mean\_sales =} \FunctionTok{mean}\NormalTok{(time}\FloatTok{.2}\NormalTok{, }\AttributeTok{na.rm =} \ConstantTok{TRUE}\NormalTok{),}
    \AttributeTok{se =} \FunctionTok{sd}\NormalTok{(time}\FloatTok{.2}\NormalTok{, }\AttributeTok{na.rm =} \ConstantTok{TRUE}\NormalTok{) }\SpecialCharTok{/} \FunctionTok{sqrt}\NormalTok{(}\FunctionTok{n}\NormalTok{()),}
    \AttributeTok{.groups =} \StringTok{\textquotesingle{}drop\textquotesingle{}}
\NormalTok{  )}

\FunctionTok{ggplot}\NormalTok{(summary\_data, }\FunctionTok{aes}\NormalTok{(}\AttributeTok{x =}\NormalTok{ biz\_focus, }\AttributeTok{y =}\NormalTok{ mean\_sales, }\AttributeTok{fill =}\NormalTok{ internet)) }\SpecialCharTok{+}
  \FunctionTok{geom\_bar}\NormalTok{(}\AttributeTok{stat =} \StringTok{"identity"}\NormalTok{, }\AttributeTok{position =} \FunctionTok{position\_dodge}\NormalTok{(}\FloatTok{0.9}\NormalTok{), }\AttributeTok{color =} \StringTok{"black"}\NormalTok{) }\SpecialCharTok{+}
  \FunctionTok{geom\_errorbar}\NormalTok{(}\FunctionTok{aes}\NormalTok{(}\AttributeTok{ymin =}\NormalTok{ mean\_sales }\SpecialCharTok{{-}}\NormalTok{ se, }\AttributeTok{ymax =}\NormalTok{ mean\_sales }\SpecialCharTok{+}\NormalTok{ se),}
                \AttributeTok{width =} \FloatTok{0.2}\NormalTok{, }\AttributeTok{position =} \FunctionTok{position\_dodge}\NormalTok{(}\FloatTok{0.9}\NormalTok{)) }\SpecialCharTok{+}
  \FunctionTok{scale\_fill\_manual}\NormalTok{(}\AttributeTok{values =} \FunctionTok{c}\NormalTok{(}\StringTok{"white"}\NormalTok{, }\StringTok{"gray50"}\NormalTok{, }\StringTok{"gray20"}\NormalTok{)) }\SpecialCharTok{+}
  \FunctionTok{labs}\NormalTok{(}\AttributeTok{x =} \StringTok{"Business Focus"}\NormalTok{, }
       \AttributeTok{y =} \StringTok{"Mean Sales Score at Time 2"}\NormalTok{,}
       \AttributeTok{fill =} \StringTok{"Internet Use Level"}\NormalTok{) }\SpecialCharTok{+}
\NormalTok{  cleanup}
\end{Highlighting}
\end{Shaded}

\pandocbounded{\includegraphics[keepaspectratio]{05-Lab-Solution_files/figure-latex/bar1-1.pdf}}

This bar chart illustrates the interaction between business focus type
and internet use level on Time 2 sales performance, with error bars
indicating the variability within each group.

\begin{enumerate}
\def\labelenumi{\arabic{enumi})}
\setcounter{enumi}{2}
\item
  Make a bar chart with two independent variables:

  \begin{enumerate}
  \def\labelenumii{\alph{enumii}.}
  \tightlist
  \item
    Time (time 1, time 2), Business focus, DV: is sales from time 1 and
    2
  \end{enumerate}
\end{enumerate}

\begin{Shaded}
\begin{Highlighting}[]
\CommentTok{\# Reshape data for time comparison}
\NormalTok{long\_data }\OtherTok{\textless{}{-}}\NormalTok{ labdata }\SpecialCharTok{\%\textgreater{}\%}
  \FunctionTok{pivot\_longer}\NormalTok{(}\AttributeTok{cols =} \FunctionTok{c}\NormalTok{(time}\FloatTok{.1}\NormalTok{, time}\FloatTok{.2}\NormalTok{), }
               \AttributeTok{names\_to =} \StringTok{"Time"}\NormalTok{, }
               \AttributeTok{values\_to =} \StringTok{"Sales"}\NormalTok{) }\SpecialCharTok{\%\textgreater{}\%}
  \FunctionTok{mutate}\NormalTok{(}\AttributeTok{Time =} \FunctionTok{factor}\NormalTok{(Time, }
                       \AttributeTok{levels =} \FunctionTok{c}\NormalTok{(}\StringTok{"time.1"}\NormalTok{, }\StringTok{"time.2"}\NormalTok{),}
                       \AttributeTok{labels =} \FunctionTok{c}\NormalTok{(}\StringTok{"Time 1"}\NormalTok{, }\StringTok{"Time 2"}\NormalTok{)))}

\CommentTok{\# Calculate means and standard errors}
\NormalTok{summary\_time }\OtherTok{\textless{}{-}}\NormalTok{ long\_data }\SpecialCharTok{\%\textgreater{}\%}
  \FunctionTok{group\_by}\NormalTok{(Time, biz\_focus) }\SpecialCharTok{\%\textgreater{}\%}
  \FunctionTok{summarise}\NormalTok{(}
    \AttributeTok{mean\_sales =} \FunctionTok{mean}\NormalTok{(Sales, }\AttributeTok{na.rm =} \ConstantTok{TRUE}\NormalTok{),}
    \AttributeTok{se =} \FunctionTok{sd}\NormalTok{(Sales, }\AttributeTok{na.rm =} \ConstantTok{TRUE}\NormalTok{) }\SpecialCharTok{/} \FunctionTok{sqrt}\NormalTok{(}\FunctionTok{n}\NormalTok{()),}
    \AttributeTok{.groups =} \StringTok{\textquotesingle{}drop\textquotesingle{}}
\NormalTok{  )}

\FunctionTok{ggplot}\NormalTok{(summary\_time, }\FunctionTok{aes}\NormalTok{(}\AttributeTok{x =}\NormalTok{ Time, }\AttributeTok{y =}\NormalTok{ mean\_sales, }\AttributeTok{fill =}\NormalTok{ biz\_focus)) }\SpecialCharTok{+}
  \FunctionTok{geom\_bar}\NormalTok{(}\AttributeTok{stat =} \StringTok{"identity"}\NormalTok{, }\AttributeTok{position =} \FunctionTok{position\_dodge}\NormalTok{(}\FloatTok{0.9}\NormalTok{), }\AttributeTok{color =} \StringTok{"black"}\NormalTok{) }\SpecialCharTok{+}
  \FunctionTok{geom\_errorbar}\NormalTok{(}\FunctionTok{aes}\NormalTok{(}\AttributeTok{ymin =}\NormalTok{ mean\_sales }\SpecialCharTok{{-}}\NormalTok{ se, }\AttributeTok{ymax =}\NormalTok{ mean\_sales }\SpecialCharTok{+}\NormalTok{ se),}
                \AttributeTok{width =} \FloatTok{0.2}\NormalTok{, }\AttributeTok{position =} \FunctionTok{position\_dodge}\NormalTok{(}\FloatTok{0.9}\NormalTok{)) }\SpecialCharTok{+}
  \FunctionTok{scale\_fill\_manual}\NormalTok{(}\AttributeTok{values =} \FunctionTok{c}\NormalTok{(}\StringTok{"white"}\NormalTok{, }\StringTok{"gray40"}\NormalTok{)) }\SpecialCharTok{+}
  \FunctionTok{labs}\NormalTok{(}\AttributeTok{x =} \StringTok{"Time Period"}\NormalTok{, }
       \AttributeTok{y =} \StringTok{"Mean Sales Score"}\NormalTok{,}
       \AttributeTok{fill =} \StringTok{"Business Focus"}\NormalTok{) }\SpecialCharTok{+}
\NormalTok{  cleanup}
\end{Highlighting}
\end{Shaded}

\pandocbounded{\includegraphics[keepaspectratio]{05-Lab-Solution_files/figure-latex/bar2-1.pdf}}

This bar chart compares sales performance across the two time periods
for traditional versus forward-thinking business focus groups, showing
how each group's performance changed over time.

\begin{enumerate}
\def\labelenumi{\arabic{enumi})}
\setcounter{enumi}{3}
\item
  Make a simple line graph:

  \begin{enumerate}
  \def\labelenumii{\alph{enumii}.}
  \tightlist
  \item
    Time (time 1, time 2), DV: is sales from time 1 and 2
  \end{enumerate}
\end{enumerate}

\begin{Shaded}
\begin{Highlighting}[]
\CommentTok{\# Calculate overall means and standard errors by time}
\NormalTok{summary\_overall }\OtherTok{\textless{}{-}}\NormalTok{ long\_data }\SpecialCharTok{\%\textgreater{}\%}
  \FunctionTok{group\_by}\NormalTok{(Time) }\SpecialCharTok{\%\textgreater{}\%}
  \FunctionTok{summarise}\NormalTok{(}
    \AttributeTok{mean\_sales =} \FunctionTok{mean}\NormalTok{(Sales, }\AttributeTok{na.rm =} \ConstantTok{TRUE}\NormalTok{),}
    \AttributeTok{se =} \FunctionTok{sd}\NormalTok{(Sales, }\AttributeTok{na.rm =} \ConstantTok{TRUE}\NormalTok{) }\SpecialCharTok{/} \FunctionTok{sqrt}\NormalTok{(}\FunctionTok{n}\NormalTok{()),}
    \AttributeTok{.groups =} \StringTok{\textquotesingle{}drop\textquotesingle{}}
\NormalTok{  )}

\FunctionTok{ggplot}\NormalTok{(summary\_overall, }\FunctionTok{aes}\NormalTok{(}\AttributeTok{x =}\NormalTok{ Time, }\AttributeTok{y =}\NormalTok{ mean\_sales, }\AttributeTok{group =} \DecValTok{1}\NormalTok{)) }\SpecialCharTok{+}
  \FunctionTok{geom\_line}\NormalTok{(}\AttributeTok{linewidth =} \FloatTok{1.2}\NormalTok{) }\SpecialCharTok{+}
  \FunctionTok{geom\_point}\NormalTok{(}\AttributeTok{size =} \DecValTok{3}\NormalTok{) }\SpecialCharTok{+}
  \FunctionTok{geom\_errorbar}\NormalTok{(}\FunctionTok{aes}\NormalTok{(}\AttributeTok{ymin =}\NormalTok{ mean\_sales }\SpecialCharTok{{-}}\NormalTok{ se, }\AttributeTok{ymax =}\NormalTok{ mean\_sales }\SpecialCharTok{+}\NormalTok{ se),}
                \AttributeTok{width =} \FloatTok{0.1}\NormalTok{) }\SpecialCharTok{+}
  \FunctionTok{labs}\NormalTok{(}\AttributeTok{x =} \StringTok{"Time Period"}\NormalTok{, }
       \AttributeTok{y =} \StringTok{"Mean Sales Score"}\NormalTok{) }\SpecialCharTok{+}
\NormalTok{  cleanup}
\end{Highlighting}
\end{Shaded}

\pandocbounded{\includegraphics[keepaspectratio]{05-Lab-Solution_files/figure-latex/line-1.pdf}}

This line graph demonstrates the overall trend in mean sales scores from
Time 1 to Time 2 across all companies, regardless of business focus or
internet use level.

\begin{enumerate}
\def\labelenumi{\arabic{enumi})}
\setcounter{enumi}{4}
\item
  Make a simple scatterplot:

  \begin{enumerate}
  \def\labelenumii{\alph{enumii}.}
  \tightlist
  \item
    Sales at Time 1, Time 2
  \end{enumerate}
\end{enumerate}

\begin{Shaded}
\begin{Highlighting}[]
\FunctionTok{ggplot}\NormalTok{(labdata, }\FunctionTok{aes}\NormalTok{(}\AttributeTok{x =}\NormalTok{ time}\FloatTok{.1}\NormalTok{, }\AttributeTok{y =}\NormalTok{ time}\FloatTok{.2}\NormalTok{)) }\SpecialCharTok{+}
  \FunctionTok{geom\_point}\NormalTok{(}\AttributeTok{size =} \DecValTok{2}\NormalTok{, }\AttributeTok{alpha =} \FloatTok{0.6}\NormalTok{) }\SpecialCharTok{+}
  \FunctionTok{labs}\NormalTok{(}\AttributeTok{x =} \StringTok{"Sales Score at Time 1"}\NormalTok{, }
       \AttributeTok{y =} \StringTok{"Sales Score at Time 2"}\NormalTok{) }\SpecialCharTok{+}
\NormalTok{  cleanup}
\end{Highlighting}
\end{Shaded}

\pandocbounded{\includegraphics[keepaspectratio]{05-Lab-Solution_files/figure-latex/scatter1-1.pdf}}

This scatterplot shows the relationship between sales scores at Time 1
and Time 2, allowing us to examine whether companies with higher initial
sales also achieved higher sales in the follow-up period.

\begin{enumerate}
\def\labelenumi{\arabic{enumi})}
\setcounter{enumi}{5}
\item
  Make a grouped scatterplot:

  \begin{enumerate}
  \def\labelenumii{\alph{enumii}.}
  \tightlist
  \item
    Sales at time 1 and 2, Business focus
  \end{enumerate}
\end{enumerate}

\begin{Shaded}
\begin{Highlighting}[]
\FunctionTok{ggplot}\NormalTok{(labdata, }\FunctionTok{aes}\NormalTok{(}\AttributeTok{x =}\NormalTok{ time}\FloatTok{.1}\NormalTok{, }\AttributeTok{y =}\NormalTok{ time}\FloatTok{.2}\NormalTok{, }\AttributeTok{shape =}\NormalTok{ biz\_focus)) }\SpecialCharTok{+}
  \FunctionTok{geom\_point}\NormalTok{(}\AttributeTok{size =} \FloatTok{2.5}\NormalTok{, }\AttributeTok{alpha =} \FloatTok{0.7}\NormalTok{) }\SpecialCharTok{+}
  \FunctionTok{scale\_shape\_manual}\NormalTok{(}\AttributeTok{values =} \FunctionTok{c}\NormalTok{(}\DecValTok{1}\NormalTok{, }\DecValTok{16}\NormalTok{)) }\SpecialCharTok{+}
  \FunctionTok{labs}\NormalTok{(}\AttributeTok{x =} \StringTok{"Sales Score at Time 1"}\NormalTok{, }
       \AttributeTok{y =} \StringTok{"Sales Score at Time 2"}\NormalTok{,}
       \AttributeTok{shape =} \StringTok{"Business Focus"}\NormalTok{) }\SpecialCharTok{+}
\NormalTok{  cleanup}
\end{Highlighting}
\end{Shaded}

\pandocbounded{\includegraphics[keepaspectratio]{05-Lab-Solution_files/figure-latex/scatter2-1.pdf}}

This grouped scatterplot examines the relationship between Time 1 and
Time 2 sales scores separately for traditional and forward-thinking
companies, revealing whether the two business focus types show different
patterns in sales performance over time.

\end{document}
